\chapter{Cronograma}
\label{chap:cronograma}

As atividades do TCC II vão seguir as etapas descritas no Capítulo~\ref{chap:metodologia}, e os seus prazos estão na Tabela~\ref{tab:cronograma}. O TCC II será feito no semestre 2027.1, que pelo calendário universitário da UFC aprovado para 2027 vai de 1º de março a 22 de julho \cite{ufc2026calendario}, e por isso a tabela vai de março a julho de 2027. O tamanho de cada etapa leva em conta a estimativa de treinos e de horas de GPU da Tabela~\ref{tab:orcamento}, que fica na ordem de 150 a 300 horas para os 272 treinos do plano principal. O treino piloto não entra nessa conta. Os treinos de classificação são os mais longos, e por isso a validação cruzada dessa tarefa ocupa quase todo o mês de maio. Os 44 treinos opcionais só serão feitos se sobrar tempo em junho e aparecem em uma linha separada da tabela, fora do total do plano principal.

\begin{table}[H]
	\caption{\label{tab:cronograma} Cronograma de atividades do período de 2027.1}%
	\IBGEtab{}{%
		\begin{tabular}{p{8.4cm}ccccc}
			\toprule
			Atividade & Mar. & Abr. & Maio & Jun. & Jul. \\
			\midrule
			Revisão do TCC I e ajustes pedidos pela banca & X & & & & \\
			Preparação das bases e pré-processamento & X & & & & \\
			Código e treino piloto (ResNet50 e ViT-B/16) & X & X & & & \\
			Busca de hiperparâmetros (108 treinos) & & X & & & \\
			Validação cruzada da segmentação (90 treinos) & & X & X & & \\
			Validação cruzada da classificação (60 treinos) & & & X & & \\
			Modelos finais (12 treinos) e testes no DDR e no Messidor-2 & & & & X & \\
			Análise estatística e medida do custo dos modelos & & & & X & X \\
			Desenvolvimento e avaliação da solução CAD (2 treinos) & & & & X & X \\
			Treinos opcionais, se houver folga (44 treinos) & & & & X & \\
			Escrita contínua & & X & X & X & X \\
			Escrita dos resultados e das conclusões & & & & X & X \\
			Apresentação do TCC II & & & & & X \\
			\bottomrule
		\end{tabular}%
	}{%
		\Fonte{elaborada pelo autor.}%
	}
\end{table}
